\section{Evaluation plan}\label{sec:plan}

The findings above are single-regime results. We commit in advance to the following study
before any threshold is declared final. It extends roadmap items T24 and T25 of the repository.
All code, seeds and results will be versioned in the repository. We will report estimates with
confidence intervals, never only point values.

\paragraph{Level A.}
\emph{Factors:} number of reviewers (50--5{,}000); items per batch (10--300); camp split
(50/50 to 90/10) and number of camps; rating noise; ratings per reviewer and reviewers per item;
penalty ratio (5--500); and three gauge variants: implicit (current), explicit
$\sum_jf_j\iota_j=0$, and the gauge-invariant decile minimum. We will also vary batch composition:
the share of divisive items and the presence of decoys.
\emph{Metrics:} the leak $\ell$; false negatives on consensus items; false positives on partisan
items; capture-cost curves; stability across seeds.
\emph{Decision rule:} adopt the variant and penalties that keep $\ell\le0.1$ and the
false-negative rate on consensus items at or below 5\% across all batch compositions. If none
does, report that bridging cannot meet both targets at that scale.

\paragraph{Level B.}
\emph{Factors:} respondents (1{,}500--12{,}000); items per batch (4--16); number of biased items
(0 to all); shift $\delta$ (0.2--1.2); class balance (0.1--0.5); uniform and non-uniform DIF;
number and reliability of anchors (10--60); ability treatment (proxy, anchors in the likelihood,
differential gap); one or two distorting axes.
\emph{Metrics:} item-level false-positive rate with \emph{zero} biased items (the primary
target); power surfaces; the threshold as a function of anchor reliability.
\emph{Decision rule:} choose the threshold that holds the item-level false-positive rate at or
below 5\% at the lowest anchor reliability the deployment will accept, with a multiplicity
correction across the batch, and report the power at that threshold. The pilot sizes follow from
the power tables, not from prose.

\paragraph{Level C.}
Simulated reviewer populations with heterogeneous skill and strategic types: truthful, herding,
contrarian, and saboteur coalitions. We will compare the ratio score with the leave-one-out
difference score, each under the logistic map and the asymmetric update.
\emph{Metrics:} drift of reports toward the crowd; rank correlation between $E_u$ and true
skill; time to detect a long-con reviewer.

\paragraph{Coordination.}
Jittered cartels, sparse histories, honest like-minded groups and chaining attacks.
\emph{Metrics:} the detection ROC, and the influence a coalition obtains within one epoch and over
many epochs.
